% Planning van één labo. Gebruikt door de syllabus van dit labo en door de syllabus voor labo 3-5.
\textbf{Doel:} een grafiektype kiezen voor een gegeven vraag, de
standaardgrafieken van \code{plotly.express} maken, extra dimensies
toevoegen met kleur, facetten en grootte, en eenvoudige interactiviteit
inbouwen.

\begin{schedule}
13:00--13:30 & Principes & De juiste grafiek voor de vraag (verdeling, vergelijking, verband, evolutie); kritiek op een misleidende grafiek. \\
13:30--14:30 & Basis van \code{plotly.express} & Histogram, boxplot, spreidingsdiagram en lijngrafiek, met titels, aslabels en hover-informatie. \\
14:30--15:15 & Meer dimensies & \code{color}, \code{facet\_col} en \code{size}; één vaste kleur per positie. \\
15:15--15:30 & Pauze & \\
15:30--16:15 & Hartslag als tijdreeks & Zonebanden met \code{add\_hrect}, en waarom we geen tweede y-as gebruiken. \\
16:15--16:45 & Interactiviteit & \code{mo.ui.dropdown}, \code{mo.ui.range\_slider}, optioneel selecties met \code{mo.ui.plotly}. \\
16:45--17:30 & Minidashboard & 3 grafieken en 1 keuzemenu waarop minstens één grafiek reageert. \\
\end{schedule}

\textbf{Oefeningen:} (1) een histogram van het gewicht; (2) Yo-Yo-niveau
per positie, met facetten per geslacht; (3) snelheid doorheen de tijd
voor één atleet; \textbf{opdracht:} het minidashboard.
